\documentclass[jou,11pt,donotrepeattitle]{apa7}
% APA 7.ª edición — versión en español. Compilar:  latexmk -pdf criterio_snr_es.tex
\usepackage[spanish,es-tabla,es-noquoting]{babel}
% Punto decimal en todo el documento: los valores deben coincidir digito a digito
% con la version en ingles, con el repositorio y con la salida del script de analisis.
\decimalpoint
\usepackage{csquotes}
\usepackage[style=apa,backend=biber]{biblatex}
\DeclareLanguageMapping{spanish}{spanish-apa}
\addbibresource{referencias.bib}
\usepackage{amsmath,amssymb}
\usepackage{booktabs}
\usepackage{graphicx}
\usepackage{siunitx}

\title{Criterios de detenci\'on por debajo del piso de se\~nal-ruido:
la longitud de la ventana, no la tolerancia, gobierna la detecci\'on de convergencia
en comparaciones de arquitecturas}
\shorttitle{Criterios de detenci\'on bajo el piso de ruido}
\leftheader{Speranza}

\author{Maximiliano Speranza}
\affiliation{Investigador independiente \\
Universidad Tecnol\'ogica Nacional, Argentina}

\authornote{\addORCIDlink{}{}
El autor es investigador independiente, estudiante de la Tecnicatura Universitaria en
Programaci\'on (TUP) de la Universidad Tecnol\'ogica Nacional (UTN), Argentina, y egresado de
su Diplomatura en Inteligencia Artificial. La correspondencia relativa a este art\'iculo debe
dirigirse a Maximiliano Speranza. Correo electr\'onico: maximiliano.speranza@gmail.com.
El criterio de detenci\'on analizado aqu\'i fue pre-registrado y congelado con un DOI
\emph{antes} de recolectar dato alguno, lo que descarta la selecci\'on post hoc del criterio
como explicaci\'on del defecto que se reporta. El autor declara no tener conflictos de inter\'es
y no haber recibido financiamiento para este trabajo. El c\'odigo, los historiales de validaci\'on
almacenados y el script de an\'alisis est\'an disponibles p\'ublicamente (v\'ease Pr\'acticas
Abiertas).}

\abstract{%
Las comparaciones entre arquitecturas de modelado eficiente de secuencias suelen detener el
entrenamiento con una regla de la forma ``la mejora de la exactitud de validaci\'on en los
\'ultimos $w$ pasos es menor que una tolerancia $\tau$'', y luego otorgan a cada arquitectura el
presupuesto resultante. Analizamos una regla de ese tipo, pre-registrada y congelada con un DOI
antes de recolectar datos, y mostramos que no pod\'ia detectar convergencia: no porque su
tolerancia estuviera mal calibrada, sino porque la decisi\'on de detenci\'on se tomaba por debajo
del piso de se\~nal-ruido de la medici\'on. En una tarea de recuperaci\'on asociativa
multiconsulta, la condici\'on de atenci\'on lineal con regla delta mejoraba a raz\'on de $0.31$
puntos de exactitud por evaluaci\'on en el r\'egimen tratado como meseta (pendiente positiva en
siete de ocho semillas; $t = 4.20$, $p = .004$), mientras que el ruido de la diferencia evaluada
era de $1.16$ puntos, lo que da una relaci\'on se\~nal-ruido de $0.27$. Bajo esa relaci\'on
ninguna tolerancia separa los dos errores de inter\'es: las reglas que detectan una meseta
tambi\'en detienen el entrenamiento durante una mejora genuina. Mostramos que un remedio
propuesto ---medias de ventana con tolerancia calibrada al ruido--- tiene una tasa de detecci\'on
fijada algebraicamente en $\Phi(2)$ con independencia de la ventana, y habr\'ia declarado
convergencia mientras el modelo a\'un mejoraba con probabilidad $.73$. Al descomponer la ganancia
alcanzable, un factor de $7.0$ proviene de alargar la ventana de comparaci\'on y solo $1.31$ de
reemplazar la diferencia de dos puntos por una prueba de tendencia. La falla depende de la
arquitectura: las condiciones que saturan exhiben un ruido tres \'ordenes de magnitud menor y
cierran de inmediato, de modo que el presupuesto compartido covar\'ia con la cantidad que se
compara. Ofrecemos una regla basada en potencia para elegir $w$, pre-registrable sin datos
piloto. Todos los an\'alisis reutilizan historiales de validaci\'on ya almacenados y no requieren
entrenamiento adicional.}

\keywords{criterio de detenci\'on, detenci\'on temprana, potencia estad\'istica,
pre-registro, atenci\'on lineal, recuperaci\'on asociativa, presupuesto de entrenamiento,
reproducibilidad}

\begin{document}
\maketitle

\section{Introducci\'on}

Las comparaciones entre arquitecturas de modelado eficiente de secuencias descansan sobre una
decisi\'on que rara vez se examina: cu\'ando detener el entrenamiento. Una pr\'actica habitual
consiste en entrenar una arquitectura de referencia hasta que se dispara una regla de
convergencia y luego otorgar a cada arquitectura competidora el mismo n\'umero de pasos
\parencite{li2019budgeted}. El procedimiento es v\'alido solo si la convergencia puede detectarse
de manera confiable. Si la detecci\'on no lo es, el presupuesto compartido hereda esa falta de
confiabilidad, y la comparaci\'on tambi\'en.

Este art\'iculo reporta un caso en el que la detecci\'on era imposible por construcci\'on, y
caracteriza por qu\'e. El escenario es una comparaci\'on pre-registrada de composiciones de
cabezas de atenci\'on sobre recuperaci\'on asociativa multiconsulta
\parencite[MQAR, por sus siglas en ingl\'es;][]{arora2023zoology}. La regla de detenci\'on fue
especificada, congelada y publicada con un DOI antes de entrenar modelo alguno. Eso importa para
el argumento que se desarrolla aqu\'i: el defecto no puede atribuirse a una regla elegida
despu\'es de ver los resultados, porque la regla no pudo haberse elegido despu\'es de verlos.

Nuestro diagn\'ostico difiere del que sugerir\'ia inspeccionar la regla por s\'i sola. La
tolerancia era, en efecto, menor que el ruido de la m\'etrica, y resulta tentador concluir que la
regla med\'ia entonces ruido en lugar de convergencia. Esa conclusi\'on no sobrevive al contacto
con los datos. El r\'egimen al que se aplic\'o la regla no era una meseta: la arquitectura de
inter\'es a\'un mejoraba, a una tasa que una ventana m\'as larga habr\'ia detectado. La regla
fall\'o porque la \emph{decisi\'on de detenci\'on} se tomaba por debajo del piso de se\~nal-ruido
de la medici\'on, y ninguna elecci\'on de tolerancia repara una decisi\'on tomada bajo ese piso.

Nuestra contribuci\'on tiene cuatro partes. Primero, mostramos que el r\'egimen tratado como
convergido conten\'ia una tendencia de mejora estad\'isticamente detectable, de modo que el
diagn\'ostico debe formularse en t\'erminos de se\~nal-ruido y no de tolerancia. Segundo,
mostramos que recalibrar la tolerancia al ruido medido ---el remedio natural--- produce una tasa
de detecci\'on fijada algebraicamente e independiente de la ventana, y compra poder de
detecci\'on al costo de detenerse durante una mejora genuina. Tercero, descomponemos la ganancia
alcanzable y encontramos que la longitud de la ventana domina a la elecci\'on del estad\'istico
por un factor de cinco. Cuarto, mostramos que el ruido que impulsa todo esto depende de la
arquitectura, lo que convierte una elecci\'on procedimental aparentemente neutra en una variable
de confusi\'on alineada con la variable bajo estudio, y ofrecemos una regla basada en potencia
para $w$ que es pre-registrable sin datos piloto.

\section{El escenario pre-registrado}

El experimento compara tres composiciones de cabezas en un transformer peque\~no
($193\,493$ par\'ametros) sobre MQAR: \texttt{softmax} (atenci\'on est\'andar), \texttt{delta}
(atenci\'on lineal con regla delta) y \texttt{mix22} (dos cabezas softmax y dos cabezas delta en
la misma capa). La atenci\'on lineal con regla delta actualiza un estado asociativo de tama\~no
fijo corrigiendo la lectura actual antes de escribir \parencite{yang2024gated}; se sabe que las
composiciones h\'ibridas recuperan capacidad de recuerdo que las variantes puramente lineales no
alcanzan \parencite{arora2023zoology,arora2024simple}.

La exactitud de validaci\'on se evalu\'o cada $500$ pasos como la media de la exactitud top-1
sobre las cargas $\{64, 96, 128\}$, con un conjunto de validaci\'on mantenido fijo entre
evaluaciones mediante una semilla aleatoria fija. La regla pre-registrada, denominada D-004,
declara convergencia en el paso $N$ cuando
%
\begin{equation}
\mathrm{exac}(N) - \mathrm{exac}(N - w) < \tau ,
\label{eq:d004}
\end{equation}
%
con $w = 500$ y $\tau = 0.005$. Como las evaluaciones est\'an separadas por $500$ pasos, la
Ecuaci\'on~\ref{eq:d004} compara dos evaluaciones \emph{consecutivas}. Una condici\'on cierra su
primera fase de entrenamiento solo cuando las ocho semillas satisfacen la
Ecuaci\'on~\ref{eq:d004}; de lo contrario el entrenamiento contin\'ua en bloques de $2\,500$
pasos hasta un tope de $10\,000$. El presupuesto com\'un para la tabla de comparaci\'on primaria
es entonces el m\'aximo de los puntos de convergencia por condici\'on, de modo que el
comportamiento de detenci\'on de una condici\'on fija el presupuesto otorgado a todas.

\section{M\'etodo}

Todos los an\'alisis reutilizan historiales de validaci\'on almacenados durante la campa\~na
original; ning\'un modelo fue reentrenado. Los datos comprenden ocho corridas de \texttt{delta},
ocho de \texttt{mix22} y tres de \texttt{softmax}, m\'as una semilla de \texttt{delta} extendida a
$10\,000$ pasos que se emplea solo donde se indica.

\subsection{Estimaci\'on del ruido y de la tendencia}

Para cada corrida aislamos el r\'egimen tard\'io de entrenamiento ---desde el paso $4\,000$ para
\texttt{delta} y $1\,000$ para las condiciones que saturan--- y ajustamos una recta por m\'inimos
cuadrados de la exactitud de validaci\'on contra el \'indice de evaluaci\'on. Estimamos el ruido
por dos m\'etodos: la desviaci\'on est\'andar de los residuos de ese ajuste, y la desviaci\'on
est\'andar de las diferencias sucesivas dividida por $\sqrt{2}$. El segundo es insesgado bajo una
tendencia lineal y sirve, por lo tanto, como control del primero. La cantidad que gobierna la
Ecuaci\'on~\ref{eq:d004} es la desviaci\'on est\'andar de la diferencia evaluada,
$\sigma_D = \sigma\sqrt{2}$.

La tendencia se contrasta por corrida mediante la pendiente $b$ de m\'inimos cuadrados frente a
$H_0\colon b = 0$, a una cola, y entre corridas mediante una prueba $t$ de una muestra sobre las
ocho pendientes.

\subsection{Probabilidad de detecci\'on}

Reportamos dos estimaciones. La estimaci\'on param\'etrica modela las observaciones tard\'ias
como $\mathrm{exac}(N) = \mu + \varepsilon$ con $\varepsilon \sim \mathcal{N}(0, \sigma^2)$, lo
que da $P = \Phi(\tau / \sigma_D)$ para una semilla y $P^{k}$ para $k$ semillas independientes.
La estimaci\'on no param\'etrica aprovecha que la Ecuaci\'on~\ref{eq:d004}, con $w$ igual al
intervalo de evaluaci\'on, \emph{es} una diferencia sucesiva: reportamos entonces la fracci\'on
observada de diferencias sucesivas menores que $\tau$, lo que no exige supuesto distribucional
alguno ni supuesto de independencia.

\section{Resultados}

\subsection{El ruido de la m\'etrica depende de la arquitectura}

La Tabla~\ref{tab:ruido} reporta el ruido tard\'io por condici\'on. Los dos estimadores coinciden
dentro del $1\%$, lo que adem\'as indica que la autocorrelaci\'on residual es despreciable
(v\'ease Control de supuestos). La cantidad decisiva es $\tau / \sigma_D$: por debajo de $1$, la
tolerancia no puede resolver el ruido de la cantidad contra la que se la compara.

\begin{table}[htbp]
\caption{Ruido tard\'io de la exactitud de validaci\'on por condici\'on}
\label{tab:ruido}
\small
\begin{tabular}{lccccc}
\toprule
Cond. & $n$ & $\sigma_{\text{resid}}$ & $\sigma_{\text{dif}}$ & $\sigma_D$ & $\tau/\sigma_D$ \\
\midrule
\texttt{delta}   & 8 & 0.00821 & 0.00829 & 0.01162 & 0.43 \\
\texttt{mix22}   & 8 & 0.00004 & 0.00004 & 0.00006 & 86.9 \\
\texttt{softmax} & 3 & 0.00004 & 0.00004 & 0.00006 & 88.8 \\
\bottomrule
\end{tabular}

\vspace{5pt}
{\small\raggedright \textit{Nota.} $\sigma_D = \sigma_{\text{resid}}\sqrt{2}$ es la desviaci\'on
est\'andar de la diferencia evaluada por la Ecuaci\'on~\ref{eq:d004}. Las estimaciones usan ocho
evaluaciones tard\'ias por corrida para \texttt{delta} y cuatro para las condiciones que saturan;
los cocientes de estas \'ultimas deben leerse como ``dos \'ordenes de magnitud o m\'as'', no como
valores precisos.
\par}
\end{table}

Las tres condiciones difieren en $\sigma$ por unos tres \'ordenes de magnitud. Una \'unica
tolerancia fija resulta as\'i altamente informativa para dos de ellas e inservible para la
tercera.

\subsection{El r\'egimen no es una meseta}

Antes de preguntar si la regla detecta convergencia, preguntamos si hab\'ia convergencia que
detectar. La Tabla~\ref{tab:tend} contrasta la tendencia en el r\'egimen al que se aplic\'o la
regla.

\begin{table}[htbp]
\caption{Prueba de tendencia en el r\'egimen tratado como convergido (pasos 4\,000--7\,500)}
\label{tab:tend}
\small
\begin{tabular}{lrrrrc}
\toprule
Semilla & $b$ & $t$ & $p$ & Mejora & D-004 \\
\midrule
0 & $+0.00259$ & 1.86 & .056 & $+1.81$ & s\'i \\
1 & $+0.00234$ & 1.72 & .068 & $+1.64$ & no \\
2 & $+0.00421$ & 2.54 & .022 & $+2.95$ & s\'i \\
3 & $+0.00334$ & 5.12 & .001 & $+2.34$ & no \\
4 & $-0.00014$ & $-0.12$ & .547 & $-0.10$ & s\'i \\
5 & $+0.00445$ & 3.58 & .006 & $+3.11$ & no \\
6 & $+0.00141$ & 1.37 & .110 & $+0.98$ & s\'i \\
7 & $+0.00680$ & 4.13 & .003 & $+4.76$ & no \\
\bottomrule
\end{tabular}

\vspace{5pt}
{\small\raggedright \textit{Nota.} $b$ es la pendiente de m\'inimos cuadrados por evaluaci\'on;
los valores $p$ son a una cola contra $H_0\colon b = 0$. ``Mejora'' es el incremento ajustado en
puntos de exactitud a lo largo de los $3\,500$ pasos de la ventana. La \'ultima columna indica
si la Ecuaci\'on~\ref{eq:d004} declar\'o convergencia en el paso $7\,500$.
\par}
\end{table}

La pendiente es positiva en siete de ocho semillas, con media $+0.00312$ por evaluaci\'on; una
prueba $t$ de una muestra sobre las ocho pendientes rechaza el cero, $t(7) = 4.20$, $p = .004$.
Extrapolar la pendiente media predice $+1.56$ puntos de exactitud entre los pasos $7\,500$ y
$10\,000$. La \'unica semilla efectivamente extendida a $10\,000$ gan\'o $+1.39$ puntos de
capacidad en la carga~96 ($0.9167 \rightarrow 0.9306$), en concordancia con la proyecci\'on.

Esto reencuadra el problema. La condici\'on \texttt{delta} no hab\'ia convergido a los $7\,500$
pasos. Llevarla al tope de pasos no fue un artefacto de una regla ruidosa: fue el resultado
correcto alcanzado por la raz\'on equivocada. No puede reproch\'arsele a la regla no haber
detectado una convergencia que no hab\'ia ocurrido, pero tampoco puede acredit\'arsele nada,
puesto que la Tabla~\ref{tab:tend} muestra que sus clasificaciones no guardan relaci\'on
confiable con la tendencia medida.

\subsection{La decisi\'on se toma por debajo del piso de se\~nal-ruido}

Ahora pueden compararse las dos cantidades de manera directa. La se\~nal disponible para la
Ecuaci\'on~\ref{eq:d004} es un intervalo de evaluaci\'on de mejora, $b = 0.00312$; el ruido de la
diferencia evaluada es $\sigma_D = 0.01162$. Su cociente es
%
\begin{equation}
\mathrm{SNR} = \frac{b}{\sigma_D} = 0.27 .
\label{eq:snr}
\end{equation}
%
Este \'unico n\'umero gobierna todo lo que sigue. Una regla de detenci\'on debe equilibrar dos
errores: seguir entrenando despu\'es de la convergencia, y detenerse cuando a\'un queda mejora.
Con $\mathrm{SNR} < 1$ las dos tasas de error no pueden separarse mediante ninguna elecci\'on de
$\tau$, porque las distribuciones que se discriminan se solapan casi por completo. La selecci\'on
de la tolerancia no es la palanca; solo desliza una frontera de decisi\'on a lo largo de un eje
sobre el cual las dos hip\'otesis no son distinguibles.

\subsection{Probabilidad de detecci\'on}

Para \texttt{delta}, el modelo param\'etrico da $P = \Phi(0.43) = .667$ por semilla y
$P^{8} = .039$ para las ocho simult\'aneamente. La estimaci\'on no param\'etrica ---la fracci\'on
observada de diferencias sucesivas menores que $\tau$ en el mismo r\'egimen, $N = 56$--- da
$P = .554$, IC 95\% $[.423, .684]$, y $P^{8} = .009$. La estimaci\'on no param\'etrica es
preferible: no supone nada y ajusta mejor a lo observado. Cuatro de ocho semillas fueron
clasificadas como convergidas; el conteo esperado es $4.4$ bajo la estimaci\'on no param\'etrica
($P(X \leq 4) = .52$) frente a $5.3$ bajo la param\'etrica ($P(X \leq 4) = .26$).

En cualquiera de los dos casos, la condici\'on ten\'ia menos de un $4\%$ de probabilidad de
cerrar alguna vez su fase por convergencia, lo que aqu\'i es el resultado correcto ---puesto que
no hab\'ia convergido---, pero habr\'ia sido igualmente cierto si hubiera convergido.

\subsection{Recalibrar la tolerancia no resuelve el problema}

El remedio natural consiste en comparar medias de ventana en lugar de puntos, y fijar la
tolerancia a partir del ruido medido:
%
\begin{equation}
\overline{\mathrm{exac}}_{\text{\'ultimas } m} - \overline{\mathrm{exac}}_{\text{previas } m}
< c\,\sigma_D^{(m)}, \qquad \sigma_D^{(m)} = \sigma\sqrt{2/m} .
\label{eq:fix}
\end{equation}
%
Con $m = 3$ y $c = 2$, la probabilidad de detecci\'on conjunta sube de $.039$ a $.832$. Esa cifra
no es evidencia de una regla mejor. Fijar la tolerancia en $c\,\sigma_D^{(m)}$ hace que la tasa de
detecci\'on por semilla sea igual a $\Phi(c)$ \emph{id\'enticamente}, para cualquier $m$ y
cualquier $\sigma$: con $c = 2$ vale $\Phi(2) = .977$ sea cual fuere la ventana. La
Tabla~\ref{tab:desc} separa los dos ingredientes sobre nuestros datos. Toda la ganancia proviene
de relajar la tolerancia, no de promediar en ventanas.

\begin{table*}[htbp]
\caption{Descomposici\'on de la ganancia en detecci\'on}
\label{tab:desc}
\begin{tabular}{lccc}
\toprule
Variante & $\tau$ & $\sigma_D$ & $P^{8}$ \\
\midrule
D-004 tal como se pre-registr\'o          & 0.0050 & 0.01162 & .039 \\
Solo ventana ($m = 3$, $\tau$ fija)       & 0.0050 & 0.00671 & .126 \\
Solo tolerancia recalibrada ($2\sigma_D$) & 0.0232 & 0.01162 & .832 \\
Ecuaci\'on~\ref{eq:fix} completa          & 0.0134 & 0.00671 & .832 \\
\bottomrule
\end{tabular}

\vspace{5pt}
{\small\raggedright \textit{Nota.} $P^{8}$ es la probabilidad de que las ocho semillas sean
declaradas convergidas simult\'aneamente estando en una meseta. Las dos \'ultimas filas coinciden
porque fijar $\tau = 2\sigma_D$ deja la tasa por semilla en $\Phi(2)$ con independencia de la
ventana.
\par}
\end{table*}

El costo de esa relajaci\'on es el error que la tasa de detecci\'on no mide. Bajo una tasa de
mejora real $b$, la probabilidad de que la Ecuaci\'on~\ref{eq:fix} declare convergencia de todos
modos es $\Phi\!\left(c - b\,m/\sigma_D^{(m)}\right)$. A la tasa de mejora efectivamente medida,
esa probabilidad es $.73$: la regla propuesta habr\'ia detenido a \texttt{delta} mientras a\'un
ganaba alrededor de un punto de exactitud por ventana. Su mejora m\'inima detectable, con una
tasa de detenci\'on espuria del $5\%$, es de $4.07$ puntos por bloque de $2\,500$ pasos: mayor
que la ganancia de $1.39$ puntos que la semilla extendida realiz\'o efectivamente en un bloque de
ese tama\~no. Una regla ajustada para detectar mesetas de manera confiable en este r\'egimen es,
necesariamente, una regla que se detiene durante el progreso real.

\subsection{La longitud de la ventana domina al estad\'istico}

Lo que s\'i eleva el piso es la ventana. Para $n$ evaluaciones equiespaciadas, la pendiente de
m\'inimos cuadrados tiene error est\'andar
%
\begin{equation}
\mathrm{EE}(b) = \sigma \sqrt{\frac{12}{n(n^{2}-1)}} ,
\label{eq:se}
\end{equation}
%
que decrece como $n^{-3/2}$, frente a $n^{-1/2}$ del promediado. En nuestros datos, con $n = 8$
evaluaciones, $\mathrm{EE}(b) = 0.00127$ y la pendiente media medida arroja $t = 2.46$: los mismos
ocho puntos que dejan a la Ecuaci\'on~\ref{eq:d004} en $\mathrm{SNR} = 0.27$ sostienen la
detecci\'on cuando se los usa conjuntamente. La ganancia total de $9.2\times$ se descompone en
$7.0\times$ por comparar puntos separados por siete evaluaciones en lugar de una, y $1.31\times$
por usar el estad\'istico de tendencia en lugar de la diferencia entre extremos. La longitud de
la ventana ---no la elecci\'on del estad\'istico ni la tolerancia--- es el t\'ermino dominante.

De aqu\'i se obtiene una regla que puede fijarse de antemano. Dada una tasa m\'inima de mejora
$b^{*}$ que el estudio se compromete a no pasar por alto, un nivel $\alpha$ a una cola y una
potencia $1 - \beta$, la ventana debe satisfacer
%
\begin{equation}
b^{*} \geq (z_{\alpha} + z_{\beta})\, \sigma \sqrt{\frac{12}{n(n^{2}-1)}} .
\label{eq:power}
\end{equation}
%
La Tabla~\ref{tab:pot} eval\'ua la Ecuaci\'on~\ref{eq:power} al nivel de ruido medido aqu\'i.
Muestra lo que el dise\~no original habr\'ia requerido: detectar una mejora de un punto de
exactitud por cada $2\,500$ pasos exige unas once evaluaciones, mientras que D-004 us\'o dos.

\begin{table}[htbp]
\caption{Mejora detectable en funci\'on de la longitud de la ventana}
\label{tab:pot}
\small
\begin{tabular}{cccc}
\toprule
$n$ & Ventana & $\mathrm{EE}(b)$ & $b^{*}$ \\
\midrule
4  & 1\,500 & 0.00367 & 4.56 \\
6  & 2\,500 & 0.00196 & 2.44 \\
8  & 3\,500 & 0.00127 & 1.58 \\
10 & 4\,500 & 0.00090 & 1.12 \\
12 & 5\,500 & 0.00069 & 0.85 \\
16 & 7\,500 & 0.00045 & 0.55 \\
20 & 9\,500 & 0.00032 & 0.40 \\
\bottomrule
\end{tabular}

\vspace{5pt}
{\small\raggedright \textit{Nota.} $n$ = cantidad de evaluaciones en la ventana; ``Ventana'' es
su extensi\'on en pasos de entrenamiento. $\sigma = 0.00821$, $\alpha = .05$ a una cola,
potencia $= .80$. $b^{*}$ es la m\'inima tasa de mejora detectable con esa potencia, en puntos de
exactitud por cada $2\,500$ pasos.
\par}
\end{table}

La Ecuaci\'on~\ref{eq:power} requiere $\sigma$, que no se conoce antes de entrenar. Dos caminos
evitan la circularidad. Cuando el conjunto de validaci\'on se remuestrea en cada evaluaci\'on, el
error de muestreo aporta una cota inferior anal\'itica,
$\sigma \geq \sqrt{p(1-p)/n_{\text{val}}}$, calculable en tiempo de dise\~no. Cuando se lo
mantiene fijo, como aqu\'i, ese t\'ermino se anula y el ruido residual es variabilidad del
optimizador, que debe medirse; basta un piloto corto sobre cualquiera de las condiciones y, a
diferencia de recalibrar una tolerancia despu\'es de ver los resultados, lo que se fija es una
ventana, no un veredicto.

\subsection{Propagaci\'on al presupuesto de entrenamiento}

Como la regla fija el presupuesto com\'un, su dependencia de la arquitectura no es neutra. Las
condiciones que saturan tienen $\tau/\sigma_D \approx 87$ y cierran en la primera oportunidad; la
condici\'on que no satura a\'un mejora y es llevada al tope. Cualquiera sea el presupuesto que se
adopte entonces, o bien resulta insuficiente para la condici\'on que sigue aprendiendo, o bien
excede con creces lo que necesitan las que saturan, y \emph{cu\'al de las dos cosas ocurre lo
deciden las condiciones que saturan}. El presupuesto otorgado covar\'ia as\'i con el desempe\~no
en la tarea: precisamente la cantidad bajo comparaci\'on.

La magnitud es relevante. El margen de decisi\'on pre-registrado para la hip\'otesis primaria era
de $2$ puntos de exactitud; la semilla extendida de $7\,500$ a $10\,000$ pasos gan\'o $1.39$
puntos de capacidad en la carga~96. Una diferencia inducida por el presupuesto del mismo orden
que el margen de decisi\'on no es una cuesti\'on de redondeo. Esta estimaci\'on descansa sobre una
\'unica semilla extendida y es indicativa m\'as que concluyente; el an\'alisis de tendencia de la
Tabla~\ref{tab:tend}, que usa las ocho, respalda la misma direcci\'on y magnitud.

\subsection{Control de supuestos}
\label{sec:assump}

\paragraph{Autocorrelaci\'on residual.} La autocorrelaci\'on de rezago 1 de los residuos sin
tendencia es negativa en siete de ocho corridas, con mediana $-0.315$, lo que invita a corregir
$\mathrm{Var}(D)$. La correcci\'on ser\'ia espuria. Con $n = 8$ puntos y una tendencia ajustada
removida, el estimador de rezago 1 est\'a fuertemente sesgado hacia abajo incluso bajo ruido
blanco: una simulaci\'on bajo ruido blanco ($20\,000$ r\'eplicas) da
$E[\hat\rho_1] = -0.287$, con intervalo del 95\% $[-0.803, +0.310]$, lo que ubica la mediana
observada en el percentil $49$ de la distribuci\'on nula. Este es el sesgo de muestra peque\~na
de los estimadores de autocorrelaci\'on \parencite{marriott1954bias}. Los datos contienen un
control independiente: si $\rho_1 = -0.315$ fuera real,
$\sigma_{\text{dif}}/\sigma_{\text{resid}}$ deber\'ia valer $\sqrt{1-\rho_1} = 1.147$; el cociente
observado es $1.010$, lo que implica $\rho_1 = -0.02$. Tratamos entonces a los residuos como
serialmente no correlacionados. Se\~nalamos el punto porque la correcci\'on tentadora corre en la
direcci\'on que favorece al argumento, y se habr\'ia reportado como respaldo suyo.

\paragraph{Origen del ruido.} El conjunto de validaci\'on se mantiene fijo entre evaluaciones
mediante una semilla fija y comprende alrededor de $37\,000$ consultas; el error de remuestreo
aportar\'ia $\sigma \approx 0.0018$ y, al estar fijo, no aporta nada. El $\sigma = 0.0082$ medido
es por lo tanto atribuible a la variabilidad del optimizador y no al error de medici\'on.
Agrandar el conjunto de validaci\'on no lo reducir\'ia; promediar en el tiempo o promediar pesos,
s\'i.

\paragraph{Elecci\'on del inicio del r\'egimen.} Variar el inicio entre los pasos $2\,500$ y
$5\,000$ deja $\tau/\sigma_D$ entre $0.39$ y $0.51$, siempre muy por debajo de la unidad, y no
altera el signo ni la significaci\'on de la tendencia.

\paragraph{\textquestiondown Sigue la exactitud de validaci\'on a la cantidad de inter\'es?} Entre las ocho
semillas de \texttt{delta}, la exactitud de validaci\'on final correlaciona con la capacidad en
la carga~96 a $r = .96$ y en la carga~128 a $r = .99$. La m\'etrica no es ciega al eje bajo
estudio; es ruidosa en relaci\'on con su propia tasa de cambio. El problema es de potencia, no de
constructo.

\section{Discusi\'on}

\subsection{Implicancias}

El costo pr\'actico de una regla de detenci\'on por debajo del piso de ruido no es en primer
lugar el c\'omputo desperdiciado, aunque en nuestra campa\~na alrededor del $80\%$ del
presupuesto proyectado era atribuible a ella. El costo serio es que el presupuesto de
entrenamiento pasa a ser funci\'on del ruido de la m\'etrica, que a su vez es funci\'on de d\'onde
se sit\'ua la arquitectura respecto de la tarea. Las arquitecturas que saturan se detienen
temprano; las que est\'an en el r\'egimen intermedio interesante son llevadas al tope. Toda
comparaci\'on trazada a trav\'es de esa asimetr\'ia queda confundida, y la confusi\'on est\'a
alineada con la variable que el estudio existe para medir.

El episodio deja adem\'as una lecci\'on sobre el pre-registro. Congelar una regla de detenci\'on
es buena pr\'actica, pero una regla congelada sin an\'alisis de potencia congela un supuesto no
examinado sobre el ruido de la m\'etrica. Nuestro pre-registro fij\'o $\tau$ y $w$ y no dijo nada
sobre $\sigma$; lo que deber\'ia haberse registrado en su lugar es la Ecuaci\'on~\ref{eq:power},
puesto que enuncia qu\'e se compromete el estudio a detectar en vez de c\'omo calcular\'a una
diferencia. Este es el an\'alogo, para las decisiones tomadas durante el entrenamiento, de
registrar un tama\~no de efecto objetivo en lugar de un umbral de $p$.

\subsection{Relaci\'on con trabajos previos}

Que la detenci\'on temprana basada en validaci\'on es inestable en reg\'imenes ruidosos est\'a
establecido \parencite{prechelt1998automatic}, como lo est\'a la necesidad de controlar el
presupuesto de entrenamiento al comparar arquitecturas \parencite{li2019budgeted}, y se ha
mostrado que el ruido de muestreo de semillas afecta la confiabilidad de los \emph{benchmarks}
sint\'eticos \parencite{chanin2026reliable}. La contribuci\'on espec\'ifica de este trabajo es la
cadena que los vincula y su cuantificaci\'on: una decisi\'on de detenci\'on tomada por debajo del
piso de se\~nal-ruido, un piso que difiere en tres \'ordenes de magnitud entre las arquitecturas
comparadas, y una regla de presupuesto que transmite esa diferencia a la comparaci\'on.
Reportamos adem\'as, como detalle de advertencia, que las dos reparaciones naturales ---recalibrar
la tolerancia al ruido medido y promediar sobre una ventana--- atacan el t\'ermino menor y dejan
intacto el dominante.

\subsection{Limitaciones}

Este es un caso \'unico documentado, no un estudio general. Involucra una familia de tareas
(MQAR), una escala de modelo ($193$K par\'ametros), una plataforma de hardware, tres condiciones
y ocho semillas por condici\'on. No afirmamos que las reglas de tolerancia fija fallen en
general, ni que alg\'un resultado publicado quede invalidado por este mecanismo. La estimaci\'on
del efecto de presupuesto en la carga~96 descansa sobre una \'unica semilla extendida; el
an\'alisis de tendencia que la corrobora usa las ocho, pero extrapola m\'as all\'a del rango
observado. La regla de potencia de la Ecuaci\'on~\ref{eq:power} supone una tendencia de mejora
aproximadamente lineal y ruido homoced\'astico a lo largo de la ventana, lo que es una
aproximaci\'on para la din\'amica tard\'ia de entrenamiento y fallar\'ia cerca de una transici\'on
abrupta. Por \'ultimo, la Ecuaci\'on~\ref{eq:power} se deriva aqu\'i y se valida solo de manera
retrospectiva; su validaci\'on prospectiva queda como trabajo futuro y es el paso natural
siguiente de esta l\'inea.

\subsection{Recomendaciones}

Sugerimos cuatro pr\'acticas econ\'omicas. Primera: antes de fijar una regla de detenci\'on,
enunciar la m\'inima tasa de mejora que el estudio se compromete a detectar, y dimensionar la
ventana a partir de ella con la Ecuaci\'on~\ref{eq:power}; reportar la potencia resultante.
Segunda: reportar $\tau/\sigma_D$ y no $\tau$ a secas, puesto que la tolerancia es ininterpretable
sin el ruido de la cantidad contra la que se la compara. Tercera: reportar el ruido por
condici\'on, ya que una regla adecuada para una arquitectura puede ser inservible para otra
dentro del mismo estudio, y esa asimetr\'ia es una variable de confusi\'on y no un inconveniente.
Cuarta: cuando se relaje una regla para elevar su tasa de detecci\'on, reportar el error
complementario ---la probabilidad de detenerse ante la m\'inima mejora que importa---, porque la
tasa de detecci\'on por s\'i sola siempre puede llevarse al valor que se desee.

\section{Conclusi\'on}

Una regla de detenci\'on compara una mejora contra una tolerancia, pero lo que determina si puede
funcionar no es ninguna de esas cantidades por separado: es el cociente entre la mejora acumulada
a lo largo de la ventana de comparaci\'on y el ruido de esa comparaci\'on. Cuando ese cociente
est\'a por debajo de la unidad, la regla informa hacia qu\'e lado cay\'o la \'ultima evaluaci\'on,
y ninguna tolerancia recupera la decisi\'on. Alargar la ventana s\'i lo hace, y a una tasa que
vuelve barata la correcci\'on. Cuando una regla as\'i fija adem\'as el presupuesto de
entrenamiento compartido de una comparaci\'on entre arquitecturas, su sensibilidad al ruido
---dependiente de la arquitectura--- convierte una elecci\'on procedimental en una variable de
confusi\'on alineada con la variable bajo estudio. Diagnosticarlo no requiri\'o nada m\'as que los
historiales de validaci\'on que ya se estaban almacenando, y corregirlo no requiere entrenamiento
adicional.

\section*{Pr\'acticas Abiertas}

El pre-registro fue congelado y publicado antes de la recolecci\'on de datos
(DOI:~10.5281/zenodo.21495252). El c\'odigo, los historiales de validaci\'on almacenados y el
script de an\'alisis que reproduce las cifras reportadas est\'an disponibles en
\url{https://github.com/SperanzaMax/telar-ligamento}; el an\'alisis corre en CPU en menos de un
minuto. Debe consignarse una excepci\'on: la \'unica semilla de \texttt{delta} extendida a
$10\,000$ pasos, citada una vez por su ganancia de capacidad de $1.39$ puntos, se perdi\'o por
una ca\'ida de sesi\'on antes de poder versionarse, y es la \'unica cantidad de este art\'iculo
que no puede recomputarse desde el repositorio. Toda otra cifra, incluido el an\'alisis de
tendencia que predice esa ganancia de manera independiente, se deriva de las ocho corridas
versionadas.

\printbibliography

\end{document}
